% !TEX root = ../../ctfp-print.tex

\lettrine[lhang=0.17]{圏}{論のほとんどの構成は、}他のより具体的な数学の分野からの結果を一般化したものだ。積、余積、モノイド、指数などは、圏論の登場以前からよく知られていた。異なる数学の分野では異なる名前で呼ばれていたかもしれない。集合論のデカルト積、順序論の交わり（meet）、論理学の連言（conjunction）――これらはすべて、圏論的な積という抽象的な概念の具体例だ。

米田の補題（Yoneda lemma）はこの点で際立っている。他の数学の分野にほとんど先例のない、圏一般についての包括的な命題だ。最も近い類似物はケイリーの定理（Cayley's theorem）――すべての群はある集合の置換群と同型である――だとも言われる。

米田の補題の舞台は、任意の圏 $\cat{C}$ と、$\cat{C}$ から $\Set$ への関手 $F$ だ。前節で、一部の $\Set$ 値関手は表現可能、すなわちHom関手と自然同型であることを見た。米田の補題は、すべての $\Set$ 値関手が自然変換を通じてHom関手から得られることを教えてくれる。さらにそれらすべての変換を明示的に列挙する。

自然変換について話したとき、自然性条件はかなり制限的になりうると述べた。ある対象における自然変換の成分を定義すると、自然性はその成分を「搬送」するのに十分強力になり、射によって接続された別の対象へと移される。始域圏と終域圏の対象間に射が多いほど、自然変換の成分を搬送するための制約が多くなる。$\Set$ は非常に射が豊富な圏だ。

米田の補題は、Hom関手と他の任意の関手 $F$ との間の自然変換は、たった一点での単一成分の値を指定するだけで完全に決まると教えてくれる！残りの自然変換は自然性条件から自動的に従う。

では、米田の補題に登場する二つの関手の間の自然性条件を確認しよう。最初の関手はHom関手だ。$\cat{C}$ の中の任意の対象 $x$ を射の集合 $\cat{C}(a, x)$ へと写す――ここで $a$ は $\cat{C}$ の中の固定された対象だ。また、$x \to y$ の射 $f$ を $\cat{C}(a, f)$ へと写すことも見た。

第二の関手は、任意の $\Set$ 値関手 $F$ だ。

これら二つの関手の間の自然変換を $\alpha$ と呼ぼう。$\Set$ の中で動いているので、自然変換の成分、例えば $\alpha_x$ や $\alpha_y$ は、集合間の普通の関数だ：
\begin{gather*}
  \alpha_x \Colon \cat{C}(a, x) \to F x \\
  \alpha_y \Colon \cat{C}(a, y) \to F y
\end{gather*}

\begin{figure}[H]
  \centering
  \includegraphics[width=0.4\textwidth]{images/yoneda1.png}
\end{figure}

\noindent
これらは単なる関数なので、特定の点での値を調べることができる。では、集合 $\cat{C}(a, x)$ における「点」とは何か？ここが重要な観察だ：集合 $\cat{C}(a, x)$ の各点は、$a$ から $x$ への射 $h$ でもある。

したがって、$\alpha$ の自然性の正方形：
\[\alpha_y \circ \cat{C}(a, f) = F f \circ \alpha_x\]
は、$h$ に作用させると、点ごとに：
\[\alpha_y (\cat{C}(a, f) h) = (F f) (\alpha_x h)\]
となる。前節で、Hom関手 $\cat{C}(a,-)$ の射 $f$ への作用は前合成として定義されることを思い出そう：
\[\cat{C}(a, f) h = f \circ h\]
これにより：
\[\alpha_y (f \circ h) = (F f) (\alpha_x h)\]
この条件がいかに強力かは、$x = a$ に特殊化することで見えてくる。

\begin{figure}[H]
  \centering
  \includegraphics[width=0.4\textwidth]{images/yoneda2.png}
\end{figure}

\noindent
その場合、$h$ は $a$ から $a$ への射になる。そのような射は少なくとも一つ存在することが知られている。$h = \id_a$ だ。代入してみよう：
\[\alpha_y f = (F f) (\alpha_a \id_a)\]
何が起きたか気づいただろうか：左辺は $\cat{C}(a, y)$ の任意の元 $f$ への $\alpha_y$ の作用だ。そしてそれは $\id_a$ における $\alpha_a$ の単一の値によって完全に決まる。任意の値を選べばそれが自然変換を生成する。$\alpha_a$ の値は集合 $F a$ の中にあるので、$F a$ の任意の点がある $\alpha$ を定義する。

逆に、$\cat{C}(a, -)$ から $F$ への任意の自然変換 $\alpha$ が与えられれば、$\id_a$ での値を評価することで $F a$ の点を得られる。

以上で米田の補題が証明された：

\begin{quote}
  $\cat{C}(a, -)$ から $F$ への自然変換と $F a$ の元の間には一対一対応が存在する。
\end{quote}
言い換えると、
\[\mathit{Nat}(\cat{C}(a, -), F) \cong F a\]
あるいは、$\cat{C}$ と $\Set$ の間の関手圏を $[\cat{C}, \Set]$ と書けば、自然変換の集合はその圏のHom集合であり、次のように書ける：
\[[\cat{C}, \Set](\cat{C}(a, -), F) \cong F a\]
この対応が実際に自然同型であることは後で説明する。

ではこの結果の直観を得ようとしよう。最も驚くべきことは、自然変換全体がたった一つの核形成サイト、すなわち $\id_a$ に割り当てた値から結晶化するということだ。その点から始まり、自然性条件に従って広がっていく。$\Set$ における $\cat{C}$ の像へと溢れ出す。では、$\cat{C}(a, -)$ の下での $\cat{C}$ の像を考えてみよう。

まず $a$ 自身の像から始めよう。Hom関手 $\cat{C}(a, -)$ の下で、$a$ は集合 $\cat{C}(a, a)$ へと写される。関手 $F$ の下では、集合 $F a$ へと写される。自然変換の成分 $\alpha_a$ は $\cat{C}(a, a)$ から $F a$ への関数だ。集合 $\cat{C}(a, a)$ の中のただ一点、射 $\id_a$ に対応する点に着目しよう。それが単なる集合の一点であることを強調するために、$p$ と呼ぼう。成分 $\alpha_a$ は $p$ を $F a$ のある点 $q$ へと写すべきだ。任意の $q$ の選択が一意な自然変換を生成することを示そう。

\begin{figure}[H]
  \centering
  \includegraphics[width=0.3\textwidth]{images/yoneda3.png}
\end{figure}

\noindent
最初の主張は、一点 $q$ の選択が関数 $\alpha_a$ の残りを一意に決定するということだ。実際、$\cat{C}(a, a)$ の中の別の点 $p'$ を、$a$ から $a$ への何らかの射 $g$ に対応するものとして取ろう。ここで米田の補題の魔法が起きる：$g$ は集合 $\cat{C}(a, a)$ の点 $p'$ として見ることができる。同時に、集合間の二つの\emph{関数}を選択する。実際、Hom関手の下では射 $g$ は関数 $\cat{C}(a, g)$ へと写され、$F$ の下では $F g$ へと写される。

\begin{figure}[H]
  \centering
  \includegraphics[width=0.4\textwidth]{images/yoneda4.png}
\end{figure}

\noindent
次に、$\cat{C}(a, g)$ の $p$（覚えているように $\id_a$ に対応する）への作用を考えよう。それは前合成として定義され、$g \circ \id_a$ は $g$ に等しく、これは点 $p'$ に対応する。したがって、射 $g$ は $p$ に作用すると $p'$（つまり $g$）を生成する関数へと写される。円が完成した！

次に $F g$ の $q$ への作用を考えよう。それは $F a$ のある点 $q'$ だ。自然性の正方形を完成させるには、$p'$ が $\alpha_a$ の下で $q'$ へと写されなければならない。任意の $p'$（任意の $g$）を選び、$\alpha_a$ の下での写像を導いた。関数 $\alpha_a$ はこうして完全に決まる。

第二の主張は、$a$ に接続されている $\cat{C}$ の任意の対象 $x$ について、$\alpha_x$ が一意に決まるということだ。推論は同様だが、今度は二つの追加集合 $\cat{C}(a, x)$ と $F x$ があり、$a$ から $x$ への射 $g$ はHom関手の下で：
\[\cat{C}(a, g) \Colon \cat{C}(a, a) \to \cat{C}(a, x)\]
に写され、$F$ の下で：
\[F g \Colon F a \to F x\]
に写される。再び、$\cat{C}(a, g)$ の $p$ への作用は前合成 $g \circ \id_a$ により与えられ、$\cat{C}(a, x)$ の点 $p'$ に対応する。自然性により $p'$ への $\alpha_x$ の作用は：
\[q' = (F g) q\]
と決まる。$p'$ は任意だったので、関数 $\alpha_x$ 全体が決まる。

\begin{figure}[H]
  \centering
  \includegraphics[width=0.4\textwidth]{images/yoneda5.png}
\end{figure}

\noindent
$\cat{C}$ の中に $a$ への接続を持たない対象がある場合はどうか？それらはすべて $\cat{C}(a, -)$ の下で単一の集合――空集合――へと写される。空集合は集合の圏における始対象であることを思い出そう。これは空集合から他の任意の集合への一意な関数が存在することを意味する。その関数を \code{absurd} と呼んだ。ここでも、自然変換の成分の選択の余地はない：\code{absurd} しかありえない。

米田の補題を理解する一つの方法は、$\Set$ 値関手の間の自然変換が単なる関数の族であり、関数は一般的に情報を失うことに気づくことだ。関数は情報を潰すかもしれないし、余域の一部だけをカバーするかもしれない。情報を失わない関数は可逆なもの――同型射――だけだ。したがって、最も構造を保存する $\Set$ 値関手は表現可能なものだ。それらはHom関手か、Hom関手と自然同型な関手だ。他の関手 $F$ はすべて、情報を失う変換を通じてHom関手から得られる。そのような変換は情報を失うだけでなく、$\Set$ における関手 $F$ の像のごく一部しかカバーしないかもしれない。

\section{HaskellにおけるYoneda}

Haskellでは、Hom関手をreader関手の姿ですでに見ている：

\src{snippet01}
reader関手は射（ここでは関数）を前合成によって写す：

\src{snippet02}
米田の補題は、reader関手が他の任意の関手へと自然に写せることを教えてくれる。

自然変換は多相関数だ。したがって関手 \code{F} が与えられれば、reader関手からそれへの写像がある：

\src{snippet03}
いつものように \code{forall} はオプションだが、自然変換のパラメトリック多相を強調するために明示的に書くのが好みだ。

米田の補題は、これらの自然変換が \code{F a} の元と一対一対応することを教えてくれる：

\begin{snipv}
forall x . (a -> x) -> F x \ensuremath{\cong} F a
\end{snipv}
この恒等式の右辺は、通常データ構造として考えるものだ。関手をコンテナとして解釈したことを覚えているだろうか？\code{F a} は \code{a} のコンテナだ。しかし左辺は関数を引数として取る多相関数だ。米田の補題は二つの表現が等価であり、同じ情報を含むことを教えてくれる。

別の言い方をすれば：次の型の多相関数を与えてくれれば：

\src{snippet04}
\code{a} のコンテナを生成できる。その秘訣は米田の補題の証明で使ったものと同じだ：\code{id} でこの関数を呼び出して \code{F a} の元を得る：

\src{snippet05}
逆も真だ：型 \code{F a} の値が与えられれば：

\src{snippet06}
正しい型の多相関数を定義できる：

\src{snippet07}
二つの表現の間を自由に行き来できる。

複数の表現を持つことの利点は、一方が他方より合成しやすいかもしれないこと、あるいは一方が特定の応用において他方より効率的かもしれないことだ。

この原理の最も単純な例示は、コンパイラ構築でよく使われるコード変換だ：継続渡しスタイル（continuation passing style、\acronym{CPS}）。それは恒等関手への米田の補題の最も単純な応用だ。\code{F} を恒等で置き換えると：

\begin{snipv}
forall r . (a -> r) -> r \ensuremath{\cong} a
\end{snipv}
この式の解釈は、任意の型 \code{a} が \code{a} の「ハンドラ」を取る関数で置き換えられるということだ。ハンドラは \code{a} を受け取り計算の残りを実行する関数――継続だ。（型 \code{r} は通常ある種のステータスコードをカプセル化する。）

このプログラミングスタイルはUI、非同期システム、並行プログラミングで非常に一般的だ。\acronym{CPS} の欠点は制御の逆転を伴うことだ。コードは生産者と消費者（ハンドラ）に分割され、容易に合成できない。非自明なウェブプログラミングをある程度やった人なら誰でも、相互作用するステートフルなハンドラによるスパゲッティコードの悪夢を知っているだろう。後で見るように、関手とモナドを賢く使うことで、\acronym{CPS} の合成的な性質をある程度回復できる。

\section{余Yoneda}

いつものように、矢印の向きを逆にすることでボーナスの構成が得られる。米田の補題は反対圏 $\cat{C}^\mathit{op}$ に適用して、反変関手の間の写像を得ることができる。

同等に、Hom関手の始域を固定する代わりに終域を固定することで余Yoneda補題（co-Yoneda lemma）を導ける。$\cat{C}$ から $\Set$ への反変Hom関手：$\cat{C}(-, a)$ を得る。反変版の米田の補題は、この関手から他の任意の反変関手 $F$ への自然変換と集合 $F a$ の元の間の一対一対応を確立する：
\[\mathit{Nat}(\cat{C}(-, a), F) \cong F a\]
余Yoneda補題のHaskell版は次のとおりだ：

\begin{snipv}
forall x . (x -> a) -> F x \ensuremath{\cong} F a
\end{snipv}
文献によっては、反変版のことをYoneda補題と呼ぶことがある。

\section{練習問題}

\begin{enumerate}
  \tightlist
  \item
        HaskellにおけるYonedaの同型を形成する二つの関数 \code{phi} と \code{psi} が互いの逆であることを示せ。

        \begin{snip}{haskell}
phi :: (forall x . (a -> x) -> F x) -> F a
phi alpha = alpha id

psi :: F a -> (forall x . (a -> x) -> F x)
psi fa h = fmap h fa
\end{snip}
  \item
        離散圏とは、恒等射以外の射を持たない圏だ。そのような圏からの関手に対して米田の補題はどのように機能するか？
  \item
        ユニットのリスト \code{{[}(){]}} はその長さ以外の情報を含まない。したがって、データ型としては整数の符号化と見なせる。空リストはゼロを符号化し、シングルトン \code{{[}(){]}}（値であって型ではない）は1を符号化し、以降同様だ。リスト関手に対する米田の補題を使って、このデータ型の別の表現を構成せよ。
\end{enumerate}

\section{参考文献}

\begin{enumerate}
  \tightlist
  \item
        \urlref{https://www.youtube.com/watch?v=TLMxHB19khE}{Catsters}のビデオ。
\end{enumerate}
